\documentclass[11pt]{article}

\usepackage[margin=1in]{geometry}
\usepackage{amsmath,amssymb,amsthm}
\usepackage{mathtools}
\usepackage{enumitem}
\usepackage{booktabs}
\usepackage{tabularx}
\usepackage[hidelinks]{hyperref}
\usepackage{xurl}
\usepackage{microtype}

\newtheorem{definition}{Definition}
\newtheorem{lemma}{Lemma}
\newtheorem{remark}{Remark}

\title{Stable Theorem-to-Publication Spines for Terse Anonymity Papers\\
\large Mathematical Roots, Worked Cuts, Release Evolution, Claim Objects, and Lineaged Releases}
\author{}
\date{Unified archive note v0.06 (2026-04-26; archive v760)}

\begin{document}
\maketitle

\begin{abstract}
The archive now has enough compact owners that later terse papers can stop repeating the same long route from theorem root to public release.
A typical question no longer needs a fresh prose chain such as ``the theorem implies the worked packet, the worked packet implies the checked answer, the evaluator notarized the run, ABOM fixed the claim object, and Release~A lineaged the public receipt.''
This note gives that cross-series route one home.
It exports a \emph{theorem-to-publication spine}: one ordered map saying where to stop when the live issue is the theorem/accountant root, the maintained worked cut, the release-evolution stage floor, the optional anti-grinding/notary bundle, the digest-bound claim object, or the lineaged release receipt.
It does not create a new theorem, packet, claim surface, or receipt.
It only owns the cross-series stop/widen discipline so that later papers may stay short without becoming ambiguous about which layer is actually carrying the claim.
The current revision adds one audit-specimen cut against the maintained worked-example bundle so that the spine is checkable as a route, not merely readable as an editorial convention.
\end{abstract}

\paragraph{Non-redundancy note.}
Synthesis~24 owns the published mathematics crosswalk; Synthesis~55 owns theorem-root-to-review routing across maintained bridge cuts; Synthesis~17 owns the concrete auditable worked instantiation; Synthesis~32 owns release-evolution stage contracts; Eval~1 owns replayable anti-grinding / notarized evaluation bundles; Anonymity~C owns the digest-bound claim object and reader-facing diff surface; Release~A owns cross-release lineage, signed receipts, and fail-closed publication policy; Synthesis~52 owns the paper-facing on-request disclosure cut; and Synthesis~74--75 own the paired answer/request terminal shortcut and its stop/widen/reopen discipline once a live question has already crossed the publication boundary into the source-only tail.\cite{series:mathcrosswalk,series:seriesspines,series:workedexample,series:releasespine,series:evalnotary,series:abomoinl,series:releaseproperty,series:challengeanswerdisclosurepackets,series:pairedterminalcitationdiscipline,series:pairedterminalcutoverrule}
This note owns only the ordered route across the publication-side layers and explicitly stops before that later disclosure/request tail.

\paragraph{Scope.}
This is not a new mechanism note, not a new request-side note, not a new evaluator packet, and not a new release receipt.
It is a route-selection note for later terse papers that need to answer one practical question honestly: \emph{what is the earliest cross-series stop that still carries the live issue on its way to public release?} Once the live issue has crossed that publication boundary into answer-side disclosure, review-tail packaging, or source-only successor closure, this note should fail closed to the narrower imported tail owners rather than growing a second ladder.

\paragraph{Exports.}
This note exports (i) the theorem-to-publication spine, (ii) an optional-notarization rule, (iii) an earliest-sufficient-publication-stop rule, (iv) a publication-boundary rule, (v) a compact stop ladder for later terse papers, (vi) one audit-specimen cut grounded in the maintained worked example, and (vii) a tiny first-reopen-owner map for stale cross-series shortcuts.

\paragraph{How later terse papers should cite this note.}
Later terse papers should cite Synthesis~76 only when the live issue is the cross-series route itself: namely which major owner layer currently carries a claim on its way from mathematics to public release, or which owner regains authority first when a stale cross-series shortcut has failed. They should instead stop directly at Synthesis~24 / the imported mechanism owner when the theorem root is live, at Synthesis~55 or the maintained worked cut in Synthesis~17 when one concrete bridge cut, publication-boundary card, or review stop is live, at Synthesis~32 when successor maintenance and release-stage locality are live, at Eval~1 when retry/grinding / verifier-notarization is live, at Anonymity~C when the claim-object identity is live, and at Release~A when public lineage / receipt gating is live. If the live issue has already crossed the publication boundary into the exact on-request audit slice, stop at Synthesis~52; if it is one concrete worked checked-answer/request-tail seam, stop at Synthesis~17; and if it is the abstract paired answer/request shortcut or its abstract stop/widen/reopen discipline, stop at Synthesis~74--75 instead of stretching this spine past public release.\cite{series:mathcrosswalk,series:seriesspines,series:workedexample,series:releasespine,series:evalnotary,series:abomoinl,series:releaseproperty,series:challengeanswerdisclosurepackets,series:pairedterminalcitationdiscipline,series:pairedterminalcutoverrule}

\section{The spine itself}

\begin{definition}[Theorem-to-publication spine]
A \emph{theorem-to-publication spine} is an ordered tuple
\[
\Pi=(M,B,U,N,C,L)
\]
of already-owned archive stops with the following roles:
\begin{itemize}[leftmargin=*]
\item $M$: the theorem/accountant root or its canonical mathematics crosswalk export;
\item $B$: the maintained bridge/worked stop that carries the theorem-root claim through concrete receipt, replay, release, and review anchors;
\item $U$: the release-evolution stop that owns successor-maintenance stage locality;
\item $N$: an optional notarization stop that owns retry / attempt-prefix / verifier-bundle truthfulness;
\item $C$: the digest-bound public claim-object stop; and
\item $L$: the lineaged public-release stop.
\end{itemize}
The tuple records route order only.
It does not replace any imported owner.
\end{definition}

\begin{definition}[Optional-notarization rule]
A theorem-to-publication spine satisfies the \emph{optional-notarization rule} if the notarization stop $N$ is cited only when the live issue is whether one replayable evaluation bundle, attempt log, spending rule, or verifier report actually makes the claim non-gameable under retries or selective disclosure.
If the live issue is only theorem content, maintained route selection, claim-object identity, or lineaged publication, then $N$ is skipped rather than inserted by habit.
\end{definition}

\begin{definition}[Earliest sufficient publication stop]
Fix one later-terse-paper question $q$.
The \emph{earliest sufficient publication stop} $s(q)$ is the leftmost entry of $\Pi$ whose owning layer already answers $q$ without silently borrowing semantics from any stop to its right.
So $s(q)$ is a cross-series minimality rule: stop as soon as the honest owner is known, but not sooner.
\end{definition}

\begin{definition}[Publication-boundary rule]
A theorem-to-publication spine satisfies the \emph{publication-boundary rule} if once the live issue has crossed from public-release truth into the exact answer-side disclosure cut, review-tail packaging, or source-only successor closure branch, the spine does not grow a seventh stop. Instead, it fails closed to the narrower imported tail owners such as Synthesis~52 for the paper-facing disclosure slice and Synthesis~74--75 for paired answer/request terminal shortcuts.
\end{definition}

\begin{remark}[Why this deserves its own note]
The archive already owns every internal layer separately.
What recurs in later terse papers is not missing machinery but repeated uncertainty about route selection.
Without one cross-series spine, small notes tend either to widen too far (re-telling theorem, worked route, evaluator card, claim object, and release receipt at once) or to stop too early (treating a claim-object digest as if it settled theorem content, or treating a lineaged receipt as if it owned the anti-grinding evidence beneath it).
The theorem-to-publication spine prevents both failure modes by giving later notes one compact stop ladder. It is not a theorem-to-disclosure/request-tail index, and the publication-boundary rule keeps it from turning into one.
\end{remark}

\begin{table}[t]
\centering
\small
\begin{tabularx}{\linewidth}{@{}lX X@{}}
\toprule
Stop & Live question & Widen only if \\
\midrule
$M$ (theorem/accountant root) & The live issue is still the mathematical statement, accountant hypothesis, or the canonical theorem-local export boundary itself. & The live question has become one concrete carried packet, review stop, release-stage floor, claim object, or public lineage object rather than the root claim. \\
$B$ (maintained bridge / worked cut) & The live issue is one concrete maintained carrier, bridge cut, packet-local ladder, or review-side stop mirrored by the worked example and abstract bridge. & The live question has become successor-maintenance staging, retry/notary truthfulness, claim-object identity, or lineaged publication rather than the maintained route. \\
$U$ (release evolution) & The live issue is what changed across releases, which release stage owns the answer, or whether a successor-maintenance shortcut is still honest. & The live question has narrowed to retry/notary truthfulness, the exact claim-object digest surface, or the final lineaged release receipt. \\
$N$ (optional notarization) & The live issue is whether one replayable evaluation bundle, attempt prefix, or verifier report made the claim non-gameable under retries. & The live question is only the claim-object identity or lineaged publication, or the anti-grinding layer is not actually under review. \\
$C$ (claim object) & The live issue is which exact claim id / hook set / digest-bound public object was signed or diffed. & The live question has become cross-release lineage / receipt gating above the claim object, or theorem / worked / notary truth below it. \\
$L$ (lineaged release) & The live issue is whether the claim object was publicly lineaged, receipted, or withheld under fail-closed publication policy. & The live issue has moved back down to theorem content, maintained route choice, release-stage locality, retry truthfulness, or claim-object identity, or has crossed the publication boundary into exact disclosure/request-tail machinery. \\
\bottomrule
\end{tabularx}
\caption{The theorem-to-publication spine as a stop ladder for later terse papers.}
\label{tab:publication-stop-ladder}
\end{table}

\begin{lemma}[Later terse papers may use the theorem-to-publication spine]
Assume a later terse paper cites only the earliest sufficient publication stop $s(q)$ for its live question $q$, and widens rightward only when the relevant widening condition in Table~\ref{tab:publication-stop-ladder} has actually become live.
Then the paper may stay short without suppressing the honest owner of the claim.
\end{lemma}

\begin{proof}
If $s(q)=M$, then the theorem/accountant root already answers the question and every later stop would only add wrappers.
If $s(q)=B$, then the maintained route already fixes the concrete carrier or review stop, while theorem-local details to the left and publication-local wrappers to the right are unnecessary.
If $s(q)=U$, then the issue is successor maintenance and release-stage locality, so the maintained bridge is now too narrow and later publication wrappers are still too wide.
If $s(q)=N$, the optional-notarization rule ensures that anti-grinding truthfulness is the live issue rather than mere claim naming.
If $s(q)=C$, then the exact digest-bound claim object is what matters, while theorem content and lineaged publication remain distinct layers.
If $s(q)=L$, then public lineage and receipt gating are the actual issue and lower layers remain imported.
In every case the cited stop is the first layer whose owner already answers $q$, so no implicit ownership inference is needed.
\end{proof}

\section{One auditable specimen cut}
The spine should be checkable against a concrete maintained cut, not merely plausible as a prose convention.
For the current Synthesis~17 artifact cut, the maintained identifiers are
\[
\texttt{claim\_id=anondht.reader\_privacy.workedexample.v0},\qquad
\texttt{release\_id=worked-example-draft-199},\qquad
\texttt{note\_version=1.99}.
\]
The following specimen uses one carried line item, \nolinkurl{fallback\_contact\_surface}, and records the earliest sufficient stop for six different questions.
It is deliberately small: every row names an already-shipped artifact and the precise kind of question it answers.

\begin{table}[t]
\centering
\small
\begin{tabularx}{\linewidth}{@{}lX X@{}}
\toprule
Stop & Concrete maintained witness & First question honestly answered there \\
\midrule
$M$ & \nolinkurl{example\_line\_item\_owner\_map.json} entry \nolinkurl{fallback\_contact\_surface}; exposure \nolinkurl{sha256:430760d378623b3a01886c203076ca3d2fa00e5cc92f270d613341942732af3e}; owner notes Synthesis~3, Synthesis~11, Synthesis~14, and Synthesis~30. & Which theorem-side semantics and deployed-surface assumptions justify the fallback contact surface. This row does not bind a public budget value or replay verdict. \\
$B$ & \nolinkurl{example\_receipt.json} line item \nolinkurl{fallback\_contact\_surface}; budget value $0.04653341$ MaxL bits per lookup; replay hook \nolinkurl{eval-fallback-cct-v1}. & Where the claim lands on the public receipt surface. This row does not by itself answer whether the checker recomputed it or whether a successor release changed it. \\
$U$ & \nolinkurl{example\_compare\_report.json} classifies the maintained comparison as \nolinkurl{replay\_changed\_surfaces}; \nolinkurl{example\_release\_spine.json} orders the release stages \nolinkurl{diff -> classify -> obligate -> close -> summarize -> support -> cite -> quote -> clause}. & Which successor-maintenance stage owns the drift question and which publication duties follow. This row does not replace the replay plan or the final public receipt. \\
$N$ & \nolinkurl{example\_verifier\_report.json} lists \nolinkurl{eval-fallback-cct-v1} with plan spec \nolinkurl{sha256:e9168dbec2cabd9115a7f3ea9aebe0a75ff910d2295e450a340716cc7e541193} and recomputes \nolinkurl{fallback\_contact\_bits\_per\_lookup}= $0.04653341$. & Whether the declared fallback-contact value was replayed by the stated checker. This stop is optional unless retry, attempt-prefix, verifier, or anti-grinding truthfulness is live. \\
$C$ & \nolinkurl{example\_abom.json} binds the claim object under release label \nolinkurl{worked-example-draft-199}; its budget summary has total $0.265904562929$ MaxL bits per lookup for \nolinkurl{Q<=500} in a 24h window, with example signing key \nolinkurl{worked-example-key}. & Which digest-bound public claim object is being named. This row does not decide whether the object passed the final release gate. \\
$L$ & \nolinkurl{example\_release\_receipt.json} binds \nolinkurl{release\_id=worked-example-draft-199}, receipt digest \nolinkurl{095e828fd09ef186757f5dd49e84769603a2b59ea77c4a8bfb0edb720bc64cc1}, ABOM digest \nolinkurl{b2d8d186498e79b5fdcebbae069f034c4b93a7897bf26aebb2a7656b75453f25}, and a policy gate with \nolinkurl{max\_total\_budget\_bits\_per\_lookup=0.3}. & Whether the claim object has a lineaged release receipt and policy-gate binding. If the question asks for omitted audit slices after this point, the route crosses the publication boundary and must reopen Synthesis~52 or the paired-tail owners rather than adding a seventh stop to $\Pi$. \\
\bottomrule
\end{tabularx}
\caption{One non-decorative specimen for the theorem-to-publication spine.  The entries are not new evidence; they are the maintained worked-example objects that make the stop ladder auditable.}
\label{tab:auditable-specimen-cut}
\end{table}

\paragraph{Audit use.}
A reviewer can now test a terse citation sentence mechanically.
If the sentence says only that the fallback contact surface is justified by its imported theorem-side owners, Table~\ref{tab:auditable-specimen-cut} stops at $M$.
If it states the numeric public line item, it must at least stop at $B$.
If it says the successor pair changed only replayed surfaces, it has moved to $U$.
If it says the fallback value was recomputed, it has moved to $N$.
If it names the signed public claim object, it has moved to $C$.
If it says the object is release-receipted under the publication gate, it has moved to $L$.
The specimen therefore turns the abstract minimality rule into a small falsifiable routing test: a citation is too early when it asserts a fact from a row to its right, and too wide when it cites rows to the right that add no live fact.

\section{A tiny first-reopen-owner card}
Once a cross-series shortcut has failed, later terse papers still need one small owner map: which layer regains authority \emph{first}? The spine should answer that directly rather than forcing every note to rediscover the route.

\begin{quote}\small
\textbf{First-reopen-owner map.}
\begin{itemize}[leftmargin=*]
\item \emph{Root changed.} If the mathematical hypothesis, exported theorem card, or accountant floor is no longer the same, reopen $M$ via the imported theorem owner or Synthesis~24; the public-side route cannot stay honest once the root statement itself has changed. If the root statement is unchanged but the real question has become which old published Mathematics head still deserves to function as the current front end, reopen Synthesis~24 as well: that narrower crosswalk now owns the root-versus-companion-versus-route-around choice for the legacy published layer.
\item \emph{Maintained route changed.} If the maintained carrier, bridge cut, packet-local ladder, or review-side stop is no longer the same, reopen $B$ via Synthesis~55 and, if needed, the concrete maintained route in Synthesis~17; the cross-series shortcut depended on a fixed maintained route that is no longer the same route.
\item \emph{Release-evolution changed.} If the base$\to$successor comparison, release-stage floor, or public-status wrapper changed, reopen $U$ via Synthesis~32; the live issue is now successor-maintenance staging rather than theorem routing or public claim naming.
\item \emph{Notary truth changed.} If the retry / attempt-prefix / verifier truthfulness question became live, reopen $N$ via Eval~1; a claim object or release receipt no longer answers whether the supporting evaluation remained non-gameable.
\item \emph{Claim object changed.} If the claim id, hook set, or digest-bound public object changed, reopen $C$ via Anonymity~C; the route can no longer skip the claim-object owner once the public object itself is different.
\item \emph{Lineage changed.} If the lineage verdict, release receipt, or fail-closed publication gate changed, reopen $L$ via Release~A; the remaining issue is public release policy and lineage rather than the lower layers.
\end{itemize}
\end{quote}

\paragraph{One tiny yes/no drill.}
Suppose a later terse paper needs only to say that one anonymous-DHT claim still rests on the same theorem/accountant root. Then it stops at $M$.
Suppose the next question instead asks which exact maintained card currently carries that claim into the archive's auditable route. Then $M$ is too early and the paper stops at $B$.
Suppose the next question asks what changed across the successor release pair. Then the honest stop is $U$.
Suppose the next question asks whether the evaluator could have retried until success. Then the honest stop is $N$.
Suppose the next question asks which exact public claim packet was signed. Then the honest stop is $C$.
Suppose the last question asks whether that packet was actually lineaged into a public release receipt. Then the honest stop is $L$. Suppose the next question instead asks what exact audit slice should travel on request once the paper-facing answer is no longer enough. Do not append a seventh stop to $\Pi$; reopen Synthesis~52 instead. Suppose the next question asks whether one paired answer/request terminal shortcut is still honest after that widening. Reopen Synthesis~74 or Synthesis~75 instead. Suppose the question instead backs up and asks whether the legacy published theorem head being cited is still the right front end at all, or whether the new series should treat that old paper only as a narrower companion or route-around historical root. Do not stretch this wider ladder to answer that editorial choice either; reopen Synthesis~24, which now owns that direct-read classification for the legacy published layer. The point is not to cite more layers; it is to stop at the first layer that already owns the answer.

\section{Worked example pointer}
The maintained worked bundle already gives the archive the concrete middle of this spine.\cite{series:workedexample}
\nolinkurl{example_series_spine.json} mirrors theorem / receipt / replay / release / review anchors for exact maintained cuts, while \nolinkurl{example_release_spine.json} and \nolinkurl{example_release_stage_walkthrough.json} expose the release-evolution stage floor for one successor-maintenance route. Synthesis~17 now also carries one compact worked publication-boundary card for the exact disclosure slice, answer-side review handoff, first bounded source-only packet, request-side terminal closure verdict, and widening back to this wider ladder; but Synthesis~74 and Synthesis~75 remain the abstract paired-tail owners once a later paper stops asking for that one concrete worked seam and instead asks for the paired sentence discipline or the abstract stop/widen/reopen rule above it. Synthesis~55 then owns the abstract theorem-root-to-review bridge across those maintained cuts, and Synthesis~32 owns the release-evolution stage vocabulary sitting inside that bridge. When the live question is whether the supporting evaluation remained non-gameable under retries, later terse papers should widen next to Eval~1; when the live question is which digest-bound public object was signed, they should widen next to Anonymity~C; and when the live question is whether that object was actually lineaged and receipted for publication, they should widen next to Release~A. If the live question has already crossed the publication boundary into the exact audit slice omitted for space, the concrete worked boundary seam, or the abstract paired answer/request tail shortcut, this spine should stop rather than grow: reopen Synthesis~52, Synthesis~17, or Synthesis~74--75 respectively.\cite{series:seriesspines,series:releasespine,series:evalnotary,series:abomoinl,series:releaseproperty,series:challengeanswerdisclosurepackets,series:workedexample,series:pairedterminalcitationdiscipline,series:pairedterminalcutoverrule}

\paragraph{Why this helps the archive.}
The archive wants short papers, but it also wants auditable route discipline.
Synthesis~76 is therefore intentionally modest.
It does not add a new packet.
It only gives the series one stable answer to a recurring editorial question: when a terse paper moves from theorem root to auditable example to release evolution to public claim and final release, which owner is actually carrying the claim at each step? This pass also makes explicit that the ladder does \emph{not} decide every question about the old published Mathematics layer: whether a legacy paper is still a true front-end root, only a narrow companion, or a route-around historical head is now delegated back down to Synthesis~24 rather than silently absorbed into the wider theorem-to-publication ladder. The answer here still ends at publication; later on-request and source-only tail machinery is real, but it belongs to narrower imported owners rather than to this spine.

\begin{thebibliography}{99}\small

\bibitem{series:mathcrosswalk}
\textit{Synthesis~24: Mathematics Crosswalk for the Anonymity Series}.
Canonical crosswalk for the previously published Mathematics-era roots and their current series entry points.

\bibitem{series:seriesspines}
\textit{Synthesis~55: Series Spines Math-to-Review Handoffs}.
Compact theorem-root-to-review bridge for later terse papers.

\bibitem{series:workedexample}
\textit{Synthesis~17: Worked Example -- Receipt Interlock}.
Maintained auditable worked route and exact bridge / packet-local / review-side cuts.

\bibitem{series:releasespine}
\textit{Synthesis~32: Stable Release Spine Contracts for Successor Maintenance}.
Release-evolution stage vocabulary and auditable stage-walkthrough discipline.

\bibitem{series:evalnotary}
\textit{Eval~1: Notarized Evaluation Certificates for Anonymous-DHT Claims}.
Replayable anti-grinding / attempt-prefix / verifier-bundle publication surface.

\bibitem{series:abomoinl}
\textit{Anonymity~C: ABOM and OINL for Anonymity Claims}.
Digest-bound public claim object and reader-facing diff surface for source-only releases.

\bibitem{series:releaseproperty}
\textit{Release~A: Release Property Ledgers and Receipts}.
Cross-release lineage, signed receipt, and fail-closed publication-policy surface.

\bibitem{series:challengeanswerdisclosurepackets}
\textit{Synthesis~52: Challenge Answer Disclosure Packets for Successor Maintenance}.
Paper-facing on-request disclosure slices and validation companions under space pressure.

\bibitem{series:pairedterminalcitationdiscipline}
\textit{Synthesis~74: Paired Terminal Citation Discipline for Stable Review Spines}.
Stable paired answer-side / request-side terminal shortcut for later terse papers.

\bibitem{series:pairedterminalcutoverrule}
\textit{Synthesis~75: Paired Terminal Cutover Rules for Stable Review Spines}.
Stop/widen/reopen discipline for the paired answer/request terminal shortcut.

\end{thebibliography}

\end{document}
